\section{Introduction}
\label{sec:introduction}

When random points are scattered on a plane and curves of varying complexity are drawn through the space, an interesting phenomenon can be observed: more complex, ``squiggly'' curves will, on average, pass closer to the random points than simpler curves, even though both are drawn without knowledge of the point locations. This geometric insight, which we formalize as the \textbf{Squiggly Line Conjecture}, may help explain why larger neural networks often learn faster from the very beginning of training.

\subsection{Core Insight}

The key intuition is that a squiggly line covers more total length as it winds through space, giving it more opportunities to accidentally pass near any given point. This is analogous to how a longer fishing line cast randomly into a pond has better chances of landing near fish than a shorter one. In the context of neural networks, this translates to randomly initialized networks with more parameters having better ``coverage'' of the solution space.

More formally, consider two functions randomly drawn from different complexity classes: a simple linear function versus a high-degree polynomial. When evaluated on random data points, the polynomial will typically achieve smaller distances to those points, not through intelligent design, but purely through geometric advantage. The polynomial's increased oscillations provide more surface area to ``accidentally'' pass near random points in the input-output space.

This geometric principle may extend to neural networks, where width corresponds to functional complexity. A randomly initialized wide network explores a richer subspace of functions, potentially positioning itself closer to target functions through coverage effects rather than optimization alone.

\subsection{The Fishing Line Analogy}

To build intuition, consider casting fishing lines into a pond with randomly distributed fish. A straight line has limited chances of intersecting fish locations. However, a curved, meandering line covers more water surface area, increasing intersection probability purely through geometry. Similarly, complex functions ``cast'' into function space have better geometric coverage of random data points.

Crucially, this advantage has limits. A line that is too tangled wastes length by crossing itself repeatedly, reducing effective coverage. This explains the saturation effect we observe: overparameterization benefits plateau once networks become so complex they waste representational capacity through redundancy.

\subsection{Connections to Deep Learning Phenomena}

Our geometric framework may help explain several aspects of deep learning:

\textbf{Lottery Ticket Hypothesis}: Some random initializations are geometrically positioned closer to target solutions, making them ``luckier'' starting points that require less movement during training. Our geometric framework offers one potential explanation for why these winning tickets exist \cite{Frankle2018TheLT, Zhou2019DeconstructingLT}.

\textbf{Overparameterization Benefits}: Large networks don't just have more capacity---they have better geometric positioning from initialization, providing optimization advantages beyond mere expressivity \cite{Belkin2018ReconcilingMM, Neyshabur2017ExploringGI, AllenZhu2018LearningAG}.

\textbf{Width vs Depth Trade-offs}: Width provides spatial coverage of function space, while depth provides hierarchical refinement. Our analysis reveals optimal architectural choices balancing these geometric factors.

\subsection{Contributions}

Our work makes the following key contributions:

\begin{enumerate}
    \item \textbf{Formal mathematical framework} for the Squiggly Line Conjecture, establishing rigorous definitions of spatial coverage and expected distances in function space
    \item \textbf{Theoretical analysis} proving scaling laws connecting parameter count to geometric coverage, with dimension-independent results
    \item \textbf{Neural network interpretation} showing how ReLU networks implement piecewise polynomial functions with controllable complexity
    \item \textbf{Comprehensive empirical validation} across multiple dimensions, architectures, and data distributions, confirming theoretical predictions
    \item \textbf{Training dynamics analysis} demonstrating how initial geometric advantage translates to faster convergence and better final performance
    \item \textbf{Quantum mechanical formulation} revealing deep connections between wave function node structure and classical coverage principles
\end{enumerate}

\subsection{Broader Implications}

This work suggests that the success of overparameterized models in machine learning may partially follow from geometric principles. Just as a meandering river covers more ground than a straight canal, complex neural networks may position themselves closer to the data they need to learn.

The implications may extend beyond machine learning to physics, biology, and computation theory. We observe that a similar $\sqrt{\Lambda}$ scaling law (where $\Lambda$ is a complexity parameter) appears across classical polynomials, neural networks, and quantum states, with potential connections to biological neural systems.

The remainder of this paper is structured as follows: Section~\ref{sec:related_work} discusses related work in neural network theory and initialization. Section~\ref{sec:preliminaries} provides the mathematical foundations and formal problem setup. Section~\ref{sec:theory} presents our main theoretical results including scaling laws and saturation analysis. Section~\ref{sec:experiments} validates our theory through comprehensive experiments across multiple settings. The appendix extends our framework to quantum mechanics, revealing universal principles governing spatial coverage across physical systems.